\section{微分熵为何不是离散熵的直接复制}
\label{sec:07-Information-Theory/01-Information-and-Entropy/03}

一份体检表里有两列连续特征：身高，单位米；等待时间，单位秒。同事按熵的公式把求和换成积分，算出两列的``信息量''，报告说等待时间那列信息更多，建议优先保留。

你把身高那列的单位从米改成毫米——数据一个字节没动，只是每个数乘了一千——重算，身高的``信息量''凭空涨了将近 10，一举反超。

没有任何信息被创造出来。出问题的是那个量本身。

\subsection{把求和换成积分，得到的不是熵}

对密度为 \(f\) 的连续变量 \(X\)，形式上照抄离散公式：

\[
h(X)=-\int f(x)\log f(x)\,dx.
\]

这个量叫微分熵。记号像，性质不像。

第一处不像：它可以是负数。取 \(X\sim\mathrm{Unif}(0,a)\)，密度是常数 \(1/a\)，积出来 \(h(X)=\log a\)。当 \(a=0.5\) 时 \(h(X)=-1\) bit。离散熵永远非负，因为概率不超过 \(1\)，\(-\log p\ge0\)；而密度不是概率，它可以大于 \(1\)——\(\mathrm{Unif}(0,0.5)\) 的密度值就是 \(2\)。这不构成什么哲学困境，只是提醒：\(f(x)\) 带着单位，\(\log f(x)\) 是在对一个有量纲的东西取对数。

第二处不像，也就是开头那个陷阱。若 \(Y=aX+b\) 且 \(a\ne0\)：

\[
h(Y)=h(X)+\log\abs{a}.
\]

米换毫米，\(a=1000\)，微分熵直接加上 \(\log_2 1000\approx 9.97\) bit。更一般地，对可逆的 \(Y=g(X)\)，由密度变换公式代入积分得

\[
h(Y)=h(X)+\E\bigl[\log\abs{\det J_g(X)}\bigr],
\]

Jacobian 行列式量的是局部体积伸缩，微分熵的变化就是这个伸缩的对数期望。换句话说，坐标变换会把纯粹的几何信息灌进这个数里。

离散熵对重命名标签完全免疫——上一节说过，公式里根本没有 \(x\)。微分熵连乘个常数都扛不住。所以``一个连续变量的微分熵是多少''这句话，不指定坐标和单位就没有答案。

\subsection{同一张图，两套刻度}

把这件事画出来最省口舌。

\begin{center}
\begin{tikzpicture}[x=1.3cm,y=1cm,font=\small,>=stealth]
\begin{scope}
  \draw[->,black!55] (-0.15,0) -- (3.4,0);
  \draw[->,black!55] (-0.15,0) -- (-0.15,2.2);
  \draw[blue!75,thick] plot[domain=0:3,samples=90] (\x,{1.7*exp(-(\x-0.9)*(\x-0.9)/0.18)});
  \draw[red!65,thick,dashed] plot[domain=0:3,samples=90] (\x,{1.7*exp(-(\x-1.7)*(\x-1.7)/0.18)});
  \foreach \t in {0,1,2,3} {\draw[black!45] (\t,0) -- (\t,-0.12);
    \node[below,font=\footnotesize] at (\t,-0.1) {\t};}
  \node[font=\footnotesize,blue!75] at (0.35,1.95) {\(f\)};
  \node[font=\footnotesize,red!70] at (2.3,1.95) {\(g\)};
  \node[font=\footnotesize,black!80] at (1.5,-0.72) {横轴以米计};
  \node[font=\footnotesize,black!70,align=center] at (1.5,-1.5)
    {\(h(f)\approx0.31\) bit\\[1pt]\(\KL(f\Vert g)\approx5.13\) bit};
\end{scope}
\begin{scope}[shift={(4.6,0)}]
  \draw[->,black!55] (-0.15,0) -- (3.4,0);
  \draw[->,black!55] (-0.15,0) -- (-0.15,2.2);
  \draw[blue!75,thick] plot[domain=0:3,samples=90] (\x,{1.7*exp(-(\x-0.9)*(\x-0.9)/0.18)});
  \draw[red!65,thick,dashed] plot[domain=0:3,samples=90] (\x,{1.7*exp(-(\x-1.7)*(\x-1.7)/0.18)});
  \foreach \t/\l in {0/0,1/100,2/200,3/300} {\draw[black!45] (\t,0) -- (\t,-0.12);
    \node[below,font=\footnotesize] at (\t,-0.1) {\l};}
  \node[font=\footnotesize,blue!75] at (0.35,1.95) {\(f\)};
  \node[font=\footnotesize,red!70] at (2.3,1.95) {\(g\)};
  \node[font=\footnotesize,black!80] at (1.5,-0.72) {横轴以厘米计};
  \node[font=\footnotesize,black!70,align=center] at (1.5,-1.5)
    {\(h(f)\approx6.95\) bit\\[1pt]\(\KL(f\Vert g)\approx5.13\) bit};
\end{scope}
\draw[black!25] (3.95,-1.95) -- (3.95,2.35);
\end{tikzpicture}

\emph{两张图的曲线一模一样，唯一的差别是横轴刻度的单位。左边报 \(0.31\) bit，右边报 \(6.95\) bit，差值恰好是 \(\log_2 100\)。而两个分布之间的 KL 散度两边都是 \(5.13\) bit——尺度伸缩在分子分母上同时出现，约掉了。要记的就是这一张：单个微分熵随刻度浮动，两个微分熵之差不动。}
\end{center}

\subsection{量化：微分熵到底是什么的残差}

微分熵不是凭空的记号，它的意义要在量化里才浮现。把连续的 \(X\) 按步长 \(\Delta\) 分格，只记录它落在哪一格，得到离散变量 \(X_\Delta\)。第 \(i\) 格的概率约为 \(f(i\Delta)\Delta\)，代进离散熵：

\[
\begin{aligned}
H(X_\Delta)
&\approx -\sum_i f(i\Delta)\Delta\,\log\bigl(f(i\Delta)\Delta\bigr)\\
&= -\sum_i f(i\Delta)\log f(i\Delta)\,\Delta\;+\;\log\frac{1}{\Delta}.
\end{aligned}
\]

第一项是 \(-\int f\log f\) 的 Riemann 和，\(\Delta\to0\) 时趋于 \(h(X)\)。于是

\[
H(X_\Delta)\approx h(X)+\log\frac{1}{\Delta}.
\]

右边第二项是分辨率成本：格子越细，可区分的状态越多，编码位数就越多，而且随 \(\Delta\to0\) 发散。真正无穷的正是这一项——把实数精确写下来要无穷多位。微分熵是扣掉这一项之后剩下的有限残差，它记录的是分布形状对编码成本的贡献，而不是编码成本本身。

这个视角一次性解释了坐标依赖。换单位时 \(h(X)\) 变了 \(\log\abs{a}\)，但量化步长 \(\Delta\) 若同步换算，\(\log(1/\Delta)\) 会变 \(-\log\abs{a}\)，两者抵消，\(H(X_\Delta)\) 一位不差。真实的编码负担当然不该因为换单位而改变。荒谬结论只在你写下 \(h(X)\) 却忘了它身后还挂着一个 \(\Delta\) 时才出现。

传感器的例子可以顺着这条式子读。把温度记录精度从 \(0.1^\circ\mathrm{C}\) 提到 \(0.01^\circ\mathrm{C}\)，\(\Delta\) 缩小十倍，每个样本的离散熵约增加 \(\log_2 10\approx3.32\) bit。多出来的不是环境变随机了，是格子多了十倍。而如果精度已经细过传感器本身的噪声，多存的小数位记的主要是噪声——文件确实变大，有效信息没有等量增加。

\subsection{有意义的是差，不是值}

既然单个微分熵随坐标漂移，凡是只用到差的量就还站得住。

互信息是最重要的一个。当三个微分熵都有定义且不出现 \(\infty-\infty\) 时，

\[
I(X;Y)=h(X)+h(Y)-h(X,Y).
\]

对 \(X\) 和 \(Y\) 各做可逆变换，\(h(X)\) 与 \(h(Y)\) 各自多出一个 Jacobian 期望项，而 \(h(X,Y)\) 多出的恰好是这两项之和，相减归零。所以互信息不认单位：身高用米还是毫米、时间用秒还是小时，它跟标签的互信息是同一个数。这也正是特征筛选该用的量。

更稳妥的写法是绕开微分熵，直接用散度定义：

\[
I(X;Y)=\KL(P_{XY}\Vert P_XP_Y).
\]

KL 散度里两个密度之比把 Jacobian 消掉了，坐标不变性是内建的；而且即使单个微分熵发散，这个式子往往仍有明确意义。条件互信息、熵功率这类经过归一化的量同样可靠。

回到开头那份体检表。要比较身高和等待时间哪一列更值得保留，能问的问题不是``哪列微分熵大''，而是``哪列跟目标变量的互信息大''。前者取决于你用什么单位记录，后者不取决于任何单位。

同一条 Jacobian 公式在深度学习里以另一副面孔出现。归一化流（normalizing flow）把一个简单分布通过可逆网络推成复杂分布，训练时对数似然里那个 \(\log\abs{\det J}\) 项，就是这里的坐标变换项。它不是技术性的补丁，它是微分熵随坐标改变这件事的账单。

\subsection{Gaussian 的微分熵与退化情形}

若 \(X\sim\mathcal N(\mu,\sigma^2)\)，

\[
h(X)=\tfrac12\log(2\pi e\sigma^2).
\]

不含 \(\mu\)：平移不改变微分熵，符合前面的加常数规律（\(a=1\) 时增量为零）。只随尺度 \(\sigma\) 变化，也符合。多元情形

\[
h(X)=\tfrac12\log\bigl((2\pi e)^d\det\Sigma\bigr),
\]

行列式刻画的是等概率椭球的体积尺度。

这里有个实践中真会碰到的边界：若 \(\Sigma\) 奇异，分布退化到一个低维子空间，相对于 \(\R^d\) 的 Lebesgue 测度根本没有密度，上式不能用。强行代入会得到 \(-\infty\)。高维特征里存在严格线性相关的列时，这不是数值问题而是定义问题，加一点抖动 \(\Sigma+\varepsilon I\) 只是把它推回可算，不等于原量存在。

\subsection{``最大熵''在两个空间里问的不是同一个问题}

离散情形下，固定有限支撑，熵最大的是均匀分布，没有别的选项。连续情形要看你固定了什么：固定区间支撑，最大微分熵的是该区间上的均匀分布；固定均值方差，是 Gaussian；什么都不固定，把变量拉伸就能让微分熵要多大有多大，问题根本没有解。

这不是同一条定理的两个版本，是两个不同的优化问题。约束的作用因此比结论的名字重要得多。下一节把约束摆到中心：不是记住哪个分布最大熵，而是看清约束怎样决定了分布的形状，以及约束选错时会错成什么样。
